% Draft replacement for Section VI (Performance Evaluation) and Table I of the SENTINEL manuscript.
% Numbers are from sentinel_sim/RESULTS.md (compact formulation, 30 seeds). Entries marked TODO-TLE are to be filled
% from results/sweeps_tle_ntn_summary.csv once that run is final; entries marked TODO-DEVICE from the device run.

\section{Performance Evaluation}
\label{sec:performance}

We evaluate SENTINEL with the open simulator released with this paper\footnote{\url{https://github.com/tanzimulalam/Quantum_SAGIN_AoII_Updated_Simulation}, commit \texttt{TODO}.}
under the parameters of Table~\ref{tab:simulation_parameters}, all of which are the values used to produce the figures.
Each configuration is run for 60{,}000 slots (10~min) and 30 independent seeds; all policies share the topology,
source trajectories, fading, and disruption realization of a seed. We report medians with interquartile ranges, because
a few seeds contain simultaneous pass gaps on both satellite planes that dominate means; means and 95\% confidence
intervals are given in the repository.

\begin{table}[t]
\centering
\caption{Simulation Parameters (values used for all figures)}
\label{tab:simulation_parameters}
\renewcommand{\arraystretch}{1.15}
\begin{tabular}{p{0.50\columnwidth} p{0.40\columnwidth}}
\hline
\textbf{Parameter} & \textbf{Value} \\
\hline
\multicolumn{2}{l}{\emph{Network}} \\
Number of UEs & $\{5,10,15,20\}$; up to 100 in Sec.~VI-C \\
BSs / satellite serving slots & 2 / 2 (topology study: 0 to 3 each) \\
Area; BS height; UE height & $20\times20$~km$^2$; 25~m; 1.5~m \\
Satellite passes & Starlink TLE snapshot, SGP4, MTSU site; minimum elevation $25^\circ$ (synthetic model: $10^\circ$, max.\ elevation $\mathcal{U}[30^\circ,90^\circ]$) \\
Domains $K$ & 2 (one BS each; satellites shared); $K=$ number of BSs in Sec.~VI-C \\
\hline
\multicolumn{2}{l}{\emph{Channel and resources}} \\
UE transmit power & 23~dBm \\
UE--BS: carrier, bandwidth, path loss & 3.5~GHz, 20~MHz, $-147.55+20\log_{10}f+26\log_{10}d+6$~dB, BS gain 15~dBi, NF 7~dB \\
UE--SAT: carrier, bandwidth, path loss & 12~GHz, 50~MHz, free space $+1$~dB, gains 30~dBi $+$ 30~dBi, NF 9~dB \\
NTN large-scale model & TR~38.811 suburban: LOS probability, clutter loss, shadow fading per pass; atmospheric $0.5$~dB$/\sin V$ \\
Bandwidth levels $\bar\beta^q_m$ ($Q_m=3$) & $\{2.5,5,10\}$~MHz, $\Delta_\beta=2.5$~MHz \\
Spectral-efficiency cap & 6 (BS), 4 (SAT) bit/s/Hz \\
\hline
\multicolumn{2}{l}{\emph{Source and AoII}} \\
Source states; transition matrix & 4; birth--death, self-transition $\mathcal{U}[0.90,0.99]$ per UE \\
Distortion; AoII budget $\Delta^{II,\max}_i$ & $|x-\hat x|$; 10 \\
Update periods ($H_i=4$); slot; packet & $\{5,10,20,50\}$ slots; 10~ms; 20~kbit \\
\hline
\multicolumn{2}{l}{\emph{Disruption}} \\
Scenario set ($L=6$) & nominal; shadowing 6~dB/20 slots and 12~dB/30 slots; handover of each SAT (20 slots); joint loss (30 slots) \\
Realized process (per 1~s interval) & shadowing event per UE--SAT link w.p.\ 0.07; handover failure per SAT w.p.\ 0.03; natural handovers at pass changes \\
Link-failure handling (all policies) & node blocked after 60 failed slots, for 6~s; SENTINEL re-solves without the link \\
\hline
\multicolumn{2}{l}{\emph{Solver}} \\
Formulation & compact dual, $\kappa=4$, $T=6$ outer iterations, $\sigma=2$ \\
Local solver & XY-QAOA, depth 1, 512 shots, COBYLA 40 iterations (domains $\le 26$ qubits); one-hot-preserving annealing otherwise \\
ADMM $\varrho_0$, tolerance, rounds & 1 (scaled to cost spread), $10^{-2}$, 20 \\
Re-solve & on every satellite visibility change and at least every 30~s \\
Quantum backend & Qiskit Aer (noiseless); FakeTorino / FakeFez noise models; TODO-DEVICE \\
\hline
\end{tabular}
\end{table}

\textbf{Baselines.} \emph{Random} transmits with probability 0.1 per slot over a uniformly chosen visible node with
equal bandwidth sharing. \emph{AoII threshold} is a state-aware policy that transmits whenever its AoII is positive, at
most once per 5 slots, on the node with the best nominal rate. \emph{Nominal SENTINEL} runs the same solver with
$L=1$. \emph{SENTINEL} is the robust schedule with $L=6$. \emph{SENTINEL + trigger} uses SENTINEL's association and
bandwidth with transmissions triggered by the observed mismatch, rate-capped at the chosen period. The exact
reference for Sec.~VI-C is an integer linear program of $\mathscr{P}_3$ solved with HiGHS.

%-----------------------------------------------------------------------
\subsection{AoII Versus Network Size}
%-----------------------------------------------------------------------

Fig.~\ref{fig:ue_sweep} shows the median average AoII with the stochastic disruption process active.
% Numbers (synthetic passes, compact formulation, 30 seeds): see RESULTS.md "Time-domain sweeps with the compact formulation".
% Random 6.8; threshold 0.138-0.145; nominal 0.269-0.282; robust 0.261-0.268; +trigger 0.145-0.149 at 5..20 UEs.
Robust SENTINEL holds its AoII constant from 5 to 20 UEs ($0.26$ to $0.27$), whereas the nominal schedule is
$2$ to $6\%$ worse and its certified worst case is $4$ to $7\%$ higher (Fig.~\ref{fig:certificate}). The robust
schedule hedges by moving 20 to 25\% of deliveries from satellites to BSs. The state-aware threshold policy has about
half the AoII of robust SENTINEL because it transmits the moment the estimate is wrong, whereas SENTINEL samples
periodically without observing the source; SENTINEL with state-aware triggering recovers this gap ($0.145$ to
$0.149$) while retaining the worst-case certificate and raising the informative-update ratio from $0.23$ to $0.96$.

%-----------------------------------------------------------------------
\subsection{Certificate Validity, Topology, and Real Passes}
%-----------------------------------------------------------------------

The realized average AoII stayed at or below the certified worst case in 86\% of orchestration epochs
(Fig.~\ref{fig:certificate}); every violation coincides with a satellite pass gap longer than any window in
$\mathscr{L}$, which the synthetic pass model produces in two of thirty seeds. With passes propagated from the Starlink
snapshot such gaps do not occur, and the dominant satellite disruptions become natural handovers and NLOS passes.
With real passes (Table~\ref{tab:tle}) the certificate holds in 98.7\% of epochs; the remaining violations are
NLOS passes whose clutter loss exceeds the 12~dB severe-shadowing scenario. Robust and nominal schedules are
indistinguishable in realized median AoII, while the robust certificate is 4 to 8\% lower and its mean AoII at 20 UEs
is 16\% lower ($0.39$ versus $0.46$): robustness buys a guaranteed worst case at essentially no average cost. The
state-aware threshold policy, given the same ACK feedback as SENTINEL, is the strongest on average AoII ($0.13$ to
$0.15$); SENTINEL with state-aware triggering is within 10 to 20\% of it while retaining the certificate and routing
21 to 27\% of deliveries to BSs. Peak AoII is 100 to 250 for every policy, set by the 20-slot handover interruptions
and the NLOS passes rather than by scheduling. The value of robustness depends on whether the network can react to
link failures: with ACK feedback disabled for every policy, median AoII rises by an order of magnitude and the ranking
inverts, robust SENTINEL becoming the best policy ($1.35$ to $1.75$) and the state-aware threshold the worst non-random
one ($1.87$ to $2.49$), because the robust schedule keeps a quarter of the UEs on BSs that satellite shadowing cannot
reach. The robust hedge is thus worth 20 to 45\% in median AoII without failure feedback and close to nothing with it.
\begin{table}[!t]
\centering
\caption{Median Average AoII with Starlink Passes and the TR~38.811 Channel (30 seeds, disruption on)}
\label{tab:tle}
\footnotesize
\begin{tabular}{|c|c|c|c|c|c|}
\hline
\textbf{UEs} & \textbf{Random} & \textbf{AoII thr.} & \textbf{Nominal} & \textbf{SENTINEL} & \textbf{+ trigger} \\
\hline
5 & 5.94 & 0.134 & 0.253 & 0.265 & 0.155 \\
10 & 6.06 & 0.142 & 0.270 & 0.269 & 0.168 \\
15 & 5.87 & 0.142 & 0.269 & 0.264 & 0.155 \\
20 & 6.05 & 0.149 & 0.276 & 0.278 & 0.175 \\
\hline
\end{tabular}
\end{table}
Fig.~\ref{fig:topology} compares ground-only, satellite-only, and integrated deployments at 10 UEs: two BSs cover
only the UEs within 5 to 7~km of a $20\times20$~km$^2$ area, so ground-only monitoring fails (median AoII 60 to 78),
two satellite planes alone bring every policy below $0.4$, and the integrated $2+2$ and $3+3$ topologies are best,
with SENTINEL routing 21 to 30\% of deliveries to BSs as a hedge.

%-----------------------------------------------------------------------
\subsection{Validation Against the Exact Minimax Optimum}
\label{sec:exact_validation}
%-----------------------------------------------------------------------

Table~\ref{tab:optimality_gap} reports the mean gap to the exact ILP optimum of $\mathscr{P}_3$ over feasible
schedules (10 seeds; 5 for the QAOA rows).
\begin{table}[!t]
\centering
\caption{Mean Optimality Gap (\%) Above the Exact ILP Optimum of $\mathscr{P}_3$}
\label{tab:optimality_gap}
\footnotesize
\resizebox{\columnwidth}{!}{%
\begin{tabular}{|l|c|c|c|c|c|}
\hline
\textbf{Setting} & \textbf{Qubits/domain} & \textbf{XY-QAOA} & \textbf{Compact ADMM} & \textbf{Epigraph ADMM} & \textbf{Greedy+LS} \\
\hline
4--8 UEs, 2+2 & 8--24 & 0.0 (15/15 optimal) & 0.0 & 38--51 & 0.0 \\
10 / 25 UEs, 2+2 & 44 / 78 & -- & 0.0 & 42.7 / 1.8 & 0.0 \\
50 UEs, 4+3 & 142 & -- & 11.5 & 38.8 & 6.0 \\
75 / 100 UEs & 159 / 205 & -- & 49.9 / 60.5 & -- & 14.4 / 11.5 \\
Severity sweep, 20 UEs & 29--66 & -- & 0.0 & 0.2--4.7 & 0.0 \\
$L=1/3/5/6$, 20 UEs & 60--77 & -- & 0.0/3.2/6.9/0.0 & 74/233/9.0/3.8 & 0/0/1.1/0 \\
\hline
\end{tabular}%
}
\end{table}
Three observations follow. First, the XY-mixer QAOA returned the exact optimum on every instance whose domains fit the
simulator (8 to 24 qubits), with 100\% of samples one-hot and capacity-feasible; the epigraph encoding of a direct
minimax QUBO, run with the same classical search, is 38 to 51\% above the optimum on the same instances and needs 133 to
193 qubits. Second, the compact distributed pipeline is exact through 25 UEs and within 12\% at 50 UEs; its limit at
75 to 100 UEs is the coordination of five shared satellites among six to eight domains, which neither additional ADMM
rounds nor smaller dual steps resolve, and where greedy association with local search remains within 11 to 14\%.
Third, candidate pruning with $\kappa=4$ never changed the exact optimum.

%-----------------------------------------------------------------------
\subsection{Quantum Resources and Device Results}
%-----------------------------------------------------------------------

Table~\ref{tab:hardware} lists the transpiled circuits for a Heron processor (depth 1, greedy warm start) and the
results of sampling them on the device noise models. % TODO-DEVICE: add the device row(s) from results/hardware/.
\begin{table}[!t]
\centering
\caption{Hardware-Sized Circuits (6 UEs, $\kappa=4$) and Sampling Results}
\label{tab:hardware}
\footnotesize
\resizebox{\columnwidth}{!}{%
\begin{tabular}{|l|c|c|c|c|c|}
\hline
\textbf{Domain} & \textbf{Qubits} & \textbf{2q gates / depth} & \textbf{Backend} & \textbf{Feasible shots} & \textbf{Best sample optimal?} \\
\hline
0 & 18 & 137 / 169 & FakeTorino (Heron) & 54.4\% & yes \\
1 & 20 & 279 / 355 & FakeTorino (Heron) & 35.5\% & yes \\
0 & 18 & 139 / 188 & FakeFez (Heron r2) & 69.6\% & yes \\
1 & 20 & 301 / 436 & FakeFez (Heron r2) & 47.6\% & yes \\
0 & 18 & TODO & TODO-DEVICE & TODO & TODO \\
\hline
\end{tabular}%
}
\end{table}
Device noise breaks the one-hot structure of a third to two thirds of the shots; the classical decoder discards them,
and the best surviving sample equals the noiseless optimum in every case, so the schedule decoded from noisy samples is
certified at the exact ILP optimum ($\eta=2.506$). Each sampling job uses 2048 to 4096 shots and a few seconds of QPU
time; all parameter optimization is performed on the simulator.

On these domains the greedy assignment is itself optimal, and the role of the QAOA layers is to concentrate a
constrained sampler: from the $W$ state, one layer raises the probability of drawing an optimal assignment from the
uniform $0.22$ to $0.28$ and two layers to $0.42$, independently of the shot count above 128; from the greedy basis
state the optimum is retained with probability $0.99$ at depth 1. The optimum is contained in the samples in every run
because the one-hot subspace of a hardware-sized domain has at most a few hundred states. Instances on which greedy
association is suboptimal (50 UEs and above, Table~\ref{tab:optimality_gap}) require 140 to 200 qubits per domain and
are beyond simulation; demonstrating a quantum advantage on them is left to hardware with that capacity.
